\section{Hydrogen nucleus (proton) as a minimal color-singlet composite (interface + matching + audit)}
\label{app:proton_hydrogen_nucleus_interface}

\paragraph{Scope and layer discipline.}
\InterfaceTag This appendix records a minimal, falsifiable \emph{structure} statement for the hydrogen nucleus in the Standard Model stable sector: the nucleus of ordinary hydrogen is a single proton.
The goal here is not to derive nonperturbative confinement or the baryon mass spectrum from the $m=6$ anchor template; rather, we supply the smallest protocol-consistent dictionary that (i) constructs the proton as a gauge-invariant composite carrier built from the already-labeled quark degrees of freedom, and (ii) makes explicit what is assumed at the matching layer and what would count as failure.
\AuditTag No confinement or Yang--Mills mass-gap claim is made; see Appendix~\ref{app:qcd_confinement_proxy_audit} and Appendix~\ref{app:qcd_confinement_pade_pole_barrier_audit} for audit-only confinement proxies and the explicit boundary statements.

\subsection{Quantum numbers and the minimal constituent picture (matching)}
\paragraph{Standard Model quantum-number dictionary (matching).}
\MatchTag We use the PDG normalization convention $Q=T_3+Y$ and the standard quark charges
\[
Q(u)=+\frac{2}{3},\qquad Q(d)=-\frac{1}{3},
\]
and baryon number assignment $B(u)=B(d)=\frac{1}{3}$ (so that three valence quarks carry $B=1$) \cite{PDGRPP2024,PeskinSchroeder1995}.
This appendix treats these as matching-layer conventions that connect the protocol labels (Section~\ref{sec:sm_labeling_closure}) to standard SM quantum numbers.

\paragraph{Minimal constituent constraint for the hydrogen nucleus (interface).}
\InterfaceTag The hydrogen nucleus is identified with the proton, i.e.\ a stable color-singlet baryon with
\[
Q=+1,\qquad B=1.
\]
Within the minimal valence-quark dictionary above, the only $(u,d)$ composition with three constituents and $Q=+1$ is
\[
uud:\qquad Q(uud)=2\cdot\frac{2}{3}+\left(-\frac{1}{3}\right)=+1,\qquad B(uud)=3\cdot\frac{1}{3}=1.
\]
The corresponding isospin partner is the neutron $udd$ with $Q=0$ and $B=1$.

\subsection{Color singlet construction (mathematics \texorpdfstring{$\to$}{->} interface)}
\paragraph{SU(3) color as the minimal non-abelian gauge factor.}
Section~\ref{sec:sm_interface} closes a three-factor gauge redundancy with a minimal non-abelian factor identified as $SU(3)$ (Proposition~\ref{prop:channel_to_gauge}).
On the physical identification layer this factor is read as the color redundancy sector.

\paragraph{Existence of a three-quark color singlet (math).}
\MathTag Let $V\cong \CC^3$ denote the defining representation of $SU(3)$.
Then the triple tensor product contains a singlet:
\[
V^{\otimes 3}\cong \mathbf{1}\ \oplus\ \mathbf{8}\ \oplus\ \mathbf{8}\ \oplus\ \mathbf{10}.
\]
Concretely, the singlet is generated by the fully antisymmetric tensor $\varepsilon_{abc}$: if $\{e_a\}_{a=1}^3$ is a basis of $V$, the vector
\[
\ket{\mathrm{color\ singlet}}
\ \propto\
\sum_{a,b,c=1}^3 \varepsilon_{abc}\, e_a\otimes e_b\otimes e_c
\]
is invariant under the diagonal action of $SU(3)$.

\paragraph{Proton as a gauge-invariant composite carrier (interface).}
\InterfaceTag A minimal proton carrier state is therefore a composite built from three quark carriers in the fundamental color representation, projected onto the $SU(3)$ singlet sector.
This implements the general manuscript stance that \emph{composites are iterates/tensorizations, not new primitives} (see Appendix~\ref{app:composite_systems_tensor_products} for the minimal tensor-product package, and Section~\ref{subsec:gauge_as_compensation} for gauge redundancy as compensation).

\subsection{Stability gates and failure points (audit)}
\paragraph{Stability gate (what is claimed here).}
\AuditTag Within the Standard Model, the proton is stable against strong and electromagnetic decay channels; its long lifetime is therefore a consistency target for any interface dictionary that treats the proton as a persistent carrier.
In this paper's discipline, the stability claim is restricted to: \emph{given the SM quantum-number dictionary and the color-singlet composite construction above, the hydrogen nucleus is a minimal stable composite carrier}.

\paragraph{Explicit boundaries.}
\AuditTag Two boundaries are kept explicit:
\begin{itemize}
\item (P1) \textbf{No mass-spectrum closure for baryons from the $m=6$ depth template.}
The closed template in Section~\ref{sec:mass_spectrum_closure} is an anchor-level (field-label) depth assignment and is not asserted to compute nonperturbative baryon masses.
\item (P2) \textbf{No proton-decay closure.}
Any extension that predicts baryon-number violation or proton decay requires an explicit protocol-to-observation bridge and is registered as out-of-scope unless supplied (Appendix~\ref{app:unification_branching_counterfactual_audit}; Section~\ref{sec:limitations_related_work}).
\end{itemize}

\paragraph{Resolution-uplift context (interface pointer).}
\InterfaceTag The resolution staircase places a nuclear binding benchmark near $m=7$ and a conservative hadronic benchmark near $m=8$ (Section~\ref{sec:falsifiability_predictions}).
In that interface reading, the hydrogen nucleus is the minimal $B=1$ composite carrier at the onset of the hadronic band, while multi-nucleon binding is associated with the intermediate ``binding'' bridge scale.

\subsection{Vendored reference targets for the proton (matching/audit)}
\label{subsec:proton_properties_miniset}
\MatchTag Table~\ref{tab:proton_properties_miniset} records a tiny vendored reference set for proton properties used only as matching-layer targets and unit conventions (PDG/CODATA-style reference quantities). It is generated deterministically from a local JSON miniset (no network).

\AuditTag \paragraph{Audit summary (proton property miniset).} \AuditTag Rows are vendored matching-layer reference targets for the hydrogen nucleus (proton). The JSON input is stored under \texttt{data/pdg\_minisets/proton\_properties\_miniset.json}; this fragment is generated by \texttt{scripts/exp\_proton\_properties\_miniset.py}.%

\begin{table}[t]
\centering
\small
\setlength{\tabcolsep}{8pt}
\renewcommand{\arraystretch}{1.15}
\begin{tabular}{llll}
\toprule
quantity & value & unit & source \\
\midrule
electric\_charge & 1 & e & PDG RPP (convention) \\
spin & 0.5 & hbar & PDG RPP (convention) \\
mass & 0.9382720813 & GeV & CODATA/PDG (representative) \\
magnetic\_moment & 2.792847345 & mu\_N & CODATA/PDG (representative) \\
charge\_radius & 0.84 & fm & PDG/CODATA (representative; radius-puzzle caveat) \\
\end{tabular}
\caption{Vendored proton reference targets (matching layer). Rows are generated by \texttt{scripts/exp\_proton\_properties\_miniset.py} from \texttt{data/pdg\_minisets/proton\_properties\_miniset.json}.}
\label{tab:proton_properties_miniset}
\end{table}

\subsection{Resolution mapping for the proton mass (matching/audit)}
\label{subsec:proton_resolution_mapping}
\MatchTag Using the paper's resolution coordinate $r(\mu)=\log(\mu/m_e)/\log\varphi$ and the staircase convention $r_{\mathrm{step}}=2\pi$ (Section~\ref{sec:falsifiability_predictions}), Table~\ref{tab:proton_resolution_mapping} reports the induced effective window $m_{\mathrm{eff}}(\mu_p)$ for the proton mass reference.

\AuditTag \paragraph{Audit summary (proton $\mu\to r\to m_{\mathrm{eff}}$ mapping).} \AuditTag We map the vendored proton mass $\mu_p$ (GeV) to the resolution coordinate $r(\mu)=\log(\mu/m_e)/\log\varphi$ using $m_e=\texttt{M\_E\_GEV}$ and $\varphi=\texttt{PHI}$, then compute $m_{\mathrm{eff}}(\mu)=6+\lfloor r(\mu)/r_{\mathrm{step}}\rfloor$ with $r_{\mathrm{step}}=2\pi$ as in Section~\ref{sec:falsifiability_predictions}. Rows are generated by \texttt{scripts/exp\_proton\_resolution\_mapping.py} from \texttt{data/pdg\_minisets/proton\_properties\_miniset.json}.%

\begin{table}[t]
\centering
\small
\setlength{\tabcolsep}{7pt}
\renewcommand{\arraystretch}{1.15}
\begin{tabular}{lrrrrrl}
\toprule
object & $\mu$ [GeV] & $r(\mu)$ & $m_{\mathrm{eff}}$ & $\mu_{\mathrm{th}}(m_{\mathrm{eff}})$ [GeV] & $\mu_{\mathrm{th}}(m_{\mathrm{eff}}+1)$ [GeV] & note \\
\midrule
proton & 0.9382720813 & 15.6177 & 8 & 0.2160907361 & 4.443694762 & hadronic band (m=8) \\%
\end{tabular}
\caption{Resolution mapping for the proton mass reference under the fixed staircase convention $r_{\mathrm{step}}=2\pi$. Rows are generated by \texttt{scripts/exp\_proton\_resolution\_mapping.py}.}
\label{tab:proton_resolution_mapping}
\end{table}
